% !TeX root = ../../Book3.tex
\section{Schwarz–Christoffel 公式}
\label{b3:sec:0703}

多边形的边界曲率集中在顶点。直边允许 Schwarz 反射，顶点内角决定映射导数的幂次。把这些局部信息合并，就得到可积分的全局公式。

\subsection{多边形区域}
\label{b3:subsec:070301}

设 \(P\) 是有界简单多边形内部，顶点按正向顺序为 \(w_1,\ldots,w_n\)，内角为 \(\alpha_j\pi\)、\(0<\alpha_j<2\)。由 Riemann 映射与 Jordan 边界同胚，可把上半平面双全纯映到 \(P\)，并在扩充实轴上标记预顶点 \(x_j\)。选取上半平面自同构，可令无穷远对应某条边的内部点，故所有 \(x_j\) 都有限。

\term[schwarz-fanshe-yuanli]{Schwarz 反射原理}[Schwarz reflection principle]先在直边上表述。

\begin{proposition}[直边反射]\label{b3:prop:07-schwarz-reflection}
若 \(g\) 在上半圆盘全纯，连续到直径且直径上取实值，则它按 \(g(\overline z)=\overline{g(z)}\) 全纯延到整个圆盘。
\end{proposition}
\begin{proof}
定义下半盘上的反射式，直径实值使拼接连续。任何跨直径的小三角形可沿距直径 \(\varepsilon\) 的平行线分割，两侧积分为零；令 \(\varepsilon\downarrow0\)，新增切边相消。Morera 定理给拼接函数全纯。对非水平直边先作仿射旋转即可。
\end{proof}
反射要求直径上确实有连续实值。只知道上半盘全纯、或仅知道几乎处处实边界值，不足以直接套用这项初等版本。

\subsection{内角}
\label{b3:subsec:070302}

在 \(w_j\) 附近用 \((w-w_j)^{1/\alpha_j}\) 把角域拉直到半平面。分支选在多边形局部内部，连续到两条边。与 \(f\) 复合后，可跨 \(x_j\) 反射。拉直后的映射在 \(x_j\) 有一个简单零点：若阶数大于1，局部标准形将在半邻域中覆盖过大的角度，违反其到半平面的单射性与两侧对应。故有
\[
 f(z)-w_j=(z-x_j)^{\alpha_j}h_j(z),\qquad h_j(x_j)\ne0,
\]
其中幂函数在上半邻域取分支，\(h_j\) 在小邻域全纯无零。于是
\[
 \frac{f''(z)}{f'(z)}=\frac{\alpha_j-1}{z-x_j}+\text{全纯项}.
\]
凸角 \(\alpha_j<1\) 使导数发散，凹角 \(\alpha_j>1\) 使导数趋于0。这都不妨碍内部导数非零，因为顶点在边界上。

\ProseParagraphBreak
角点指数可以直接在扇形中核验。令 \(W_\alpha=\{0<\arg w<\alpha\pi\}\)，在上半平面取 \(w=z^\alpha\)。两条实半轴分别映到扇形两边；正实侧的导数为 \(\alpha x^{\alpha-1}\)。当 \(\alpha=1/2\)，直角的局部映射是平方根，导数发散如 \(x^{-1/2}\)，但函数本身连续且差值为 \(O(x^{1/2})\)。当 \(\alpha=3/2\)，凹角的局部差值为 \(O(x^{3/2})\)，导数趋于0。把角点拉直后，导数又恢复有限非零。这说明坐标必须随几何选择，不能把某个坐标中的导数大小作为不变的边界奇性。

\ProseParagraphBreak
若角域只在有限邻域内成立，幂次必须在那一邻域的内部选定。不能在环绕顶点的整个穿孔平面上不加分支约定地写 \((w-w_j)^{1/\alpha_j}\)。在实际多边形中，两条邻边夹出的内部角恰提供所需的连续辐角范围，因而局部幂次合法。余下无零因子可在小盘取对数，其 \(\alpha_j\) 次幂吸收到 \(h_j\) 中，得到正文的角点标准形。

\subsection{公式}
\label{b3:subsec:070303}

\term[schwarz-christoffel-gongshi]{Schwarz--Christoffel 公式}[Schwarz--Christoffel formula]把上述角点资料合并成映射导数。

\begin{theorem}[Schwarz--Christoffel 公式]\label{b3:thm:07-schwarz-christoffel}
在上述所有预顶点有限的归一化中，存在 \(A\in\mathbb C\)、\(C\ne0\)，使
\[
 f(z)=A+C\int_{z_0}^{z}\prod_{j=1}^n(\zeta-x_j)^{\alpha_j-1}\,\mathrm{d}\zeta.
\]
各幂在上半平面连续取分支，积分与内部路径无关。
\end{theorem}
\begin{proof}
在每个非顶点实区间，\(f'\) 的辐角是所对应边的固定方向，所以 \(f''/f'\) 取实值，并可反射到下半平面。顶点处前一小节给简单极点及留数 \(\alpha_j-1\)。无穷远对应边内点，使用坐标 \(t=1/z\) 的反射与局部单射性得 \(f'(z)=c z^{-2}(1+O(1/z))\)，故 \(f''/f'=-2/z+O(1/z^2)\)。多边形角和给 \(\sum_j(\alpha_j-1)=-2\)。减去 \(\sum_j(\alpha_j-1)/(z-x_j)\) 后，差是球面上无极点、在无穷远趋于0的函数，故由 Liouville 为0。积分对数导数，得到 \(f'=C\prod_j(z-x_j)^{\alpha_j-1}\)。上半平面单连通且乘积全纯，再取原函数得公式。
\end{proof}
这项证明从已存在的多边形映射推出公式。任意指定角度和预顶点的积分不自动是给定多边形的单叶映射；边长、循环次序和简单性仍须核对。

\subsection{参数确定}
\label{b3:subsec:070304}

公式的常数 \(A\) 负责平移，\(C\) 负责旋转和缩放，预顶点的相对位置负责多边形的形状。上半平面自同构有三个实自由度，可固定三个预顶点；其余 \(n-3\) 个实参数通常须由边长比确定。边长条件写为
\[
 |w_{j+1}-w_j|=|C|\left|\int_{x_j}^{x_{j+1}}\prod_\ell(t-x_\ell)^{\alpha_\ell-1}\,\mathrm{d}t\right|.
\]
端点奇性可积，因为 \(\alpha_j-1>-1\)。相邻区间上的常辐角自动给出转角，积分模给边长。

\ProseParagraphBreak
预顶点可以很拥挤。长而细的多边形在圆盘或半平面参数中可能对应极靠近的预顶点，这称为数值计算中的拥挤现象。它不违反存在唯一性，但会使直接解边长方程不稳定。检查公式时，尺度自由度应先消除；否则把整体放大产生的等价参数当成多个不同解，会误判非唯一性。

\ProseParagraphBreak
参数数目也可以按边界资料理解。简单 \(n\) 边形的边方向由角度和整体旋转决定，边长还满足两个实闭合条件。缩放再消去一个自由度，留下 \(n-3\) 个独立的形状参数，恰与固定三个预顶点后的数目一致。这里只是维数计数，不能据此证明非线性边长方程在每一点有唯一解；对特殊对称形状、退化形状和数值分支，仍要检查映射对应的几何条件。

\ProseParagraphBreak
对于正三角形，可利用三个相等内角 \(\alpha_j=1/3\) 和半平面自同构把预顶点固定到三个选定实点。公式导数各有 \(-2/3\) 次端点奇性，边积分仍收敛，因为 \(\int_0^\varepsilon t^{-2/3}\,\mathrm{d}t=3\varepsilon^{1/3}\)。对于矩形，各内角为 \(1/2\)，四个预顶点在固定三个之后还剩一个实参数；这就是后一节椭圆积分中的 \(k\)。三角形无需额外形状参数，四边形需要一个，这是公式自由度的具体表现。
